\subsection{正测度的像的性质}

命题 4. --- 设 T, T', T'' 是三个局部紧空间，μ 是 T 上的一个正测度，π 是一个由 T 到 T' 中的 μ-可测映射，π' 是一个由 T' 到 T'' 中的映射，且 π'' = π' ∘ π。

\begin{enumerate}
\def\labelenumi{\alph{enumi})}
\item
  设 \(\pi\) 是 \(\mu\) -逆紧的，并令 \(\mu' = \pi(\mu)\) 。为使 \(\pi'\) 是 \(\mu'\) -逆紧的，必须且只须 \(\pi''\) 是 \(\mu\) -逆紧的，此时 \(\pi''(\mu) = \pi'(\pi(\mu))\) （“测度的像的可迁性”）。
\item
  设 \(\pi'\) 是连续的，且 \(\pi''\) 是 \(\mu\) -逆紧的；则 \(\pi\) 是 \(\mu\) -逆紧的，\(\pi'\) 是 \(\pi(\mu)\) -逆紧的，且 \(\pi''(\mu) = \pi'(\pi(\mu))\) 。
\end{enumerate}

在 a) 的假设下，为使 \(\pi''\) 是 \(\mu\) -可测的，必须且只须 \(\pi'\) 是 \(\mu'\) -可测的（No.~2, 命题 3）。另一方面，若 K 是 \(T''\) 的一个紧子集，则 \(\pi''(K) = \pi^{-1}(\pi'(K))\) ；为使 \(\pi''(K)\) 是本质 \(\mu\) -可积的，根据定理 1 的推论，必须且只须 \(\pi'(K)\) 是本质 \(\mu'\) -可积的。最后，若 \(\pi''\) 是 \(\mu\) -逆紧的，则置 \(\mu'' = \pi''(\mu)\) ，我们对每个函数 \(f \in \mathcal{K}(T'')\) 有

\[
\begin{array}{r l} \int f (t ^ {\prime \prime}) d \mu^ {\prime \prime} (t ^ {\prime \prime}) & = \int f (\pi^ {\prime \prime} (t)) d \mu (t) \\ & = \int f (\pi^ {\prime} (\pi (t))) d \mu (t) = \int f (\pi^ {\prime} (t ^ {\prime})) d \mu^ {\prime} (t ^ {\prime}) \end{array}
\]

这是依 No.~2 的定理 1 得到的，这就完成了 a) 的证明。

在 b) 的假设下，设 \(K'\) 是 \(T'\) 的一个紧子集。则 \(K'' = \pi'(K')\) 是紧的，因此 \(\overline{\pi''}(K'')\) 是本质 \(\mu\) -可积的，从而 \(\overline{\pi}^{1}(K') \subset \overline{\pi''}(K'')\) 是本质 \(\mu\) -可积的（Ch. IV, §5, No.~5, 命题 7），于是 \(\pi\) 是 \(\mu\) -逆紧的。然后应用陈述的 a) 即完成证明。

推论. --- 设 T 和 T' 是两个局部紧空间，\(\mu\) 是 T 上的一个正测度，\(\pi\) 是一个由 T 到 T' 上的双射映射，\(\pi^{-1}\) 是逆映射。设 \(\pi\) 是 \(\mu\) -逆紧的，并令 \(\mu' = \pi(\mu)\) 。则 \(\pi^{-1}\) 是 \(\mu'\) -逆紧的，且 \(\pi^{-1}(\pi(\mu)) = \mu\) 。

命题 5. --- 设 T 和 X 是两个局部紧空间，\(\mu\) 是 T 上的一个正测度，\(\pi\) 是一个由 T 到 X 中的 \(\mu\) -逆紧映射，g 是定义在 X 上的有限数值函数 \(\geqslant 0\) ，且使得 \(g \circ \pi\) 对 \(\mu\) 局部可积。为使 g 对 \(\pi(\mu)\) 局部可积，必须且只须 \(\pi\) 对测度 \((g \circ \pi) \cdot \mu\) 是逆紧的，此时

\[
\pi \big ((g \circ \pi) \cdot \mu \big) = g \cdot \pi (\mu).\tag{4}
\]

置 \(\nu = \pi(\mu)\) 。为使 \(g\) 是局部 \(\nu\) -可积的，必须且只须 \(gf\) 是 \(\nu\) -可积的，对每个函数 \(f \in \mathcal{K}(X)\) 而言；由于 \(gf\) 具有紧支集，说 \(gf\) 是本质 \(\nu\) -可积的是同一回事，而这等价于说 \((g \circ \pi)(f \circ \pi)\) 是本质 \(\mu\) -可积的（定理 1）。但是，依 §5, No.~3 的定理 1，这表明 \(f \circ \pi\) 对 \(\rho = (g \circ \pi) \cdot \mu\) 是本质可积的，而依定义，这就是说 \(\pi\) 是 \(\rho\) -逆紧的（因为 \(\pi\) 显然是 \(\rho\) -可测的）。此外，

\[
\int f g d \nu = \int f (\pi (t)) g (\pi (t)) d \mu (t) = \int f (\pi (t)) d \rho (t)
\]

（No.~2, 定理 1 及 §5, 定理 1），这就证明了关系式 (4)。

命题 6. --- 设 T 和 X 是两个局部紧空间，\((\lambda_{\alpha})_{\alpha\in\mathbb{A}}\) 是 T 上的一个正测度族，按关系 \(\leqslant\) 有向，且在 \(\mathcal{M}(\mathrm{T})\) 中有一个上确界 \(\mu\) 。为使一个由 T 到 X 中的映射 \(\pi\) 是 \(\mu\) -逆紧的，必须且只须它是 \(\lambda_{\alpha}\) -逆紧的，对一切 \(\alpha\in A\) ，且族 \(\left(\pi(\lambda_{\alpha})\right)_{\alpha\in\mathbb{A}}\) 在 \(\mathcal{M}(\mathrm{X})\) 中有上界。此时，

\[
\pi (\mu) = \sup _ {\alpha} \pi (\lambda_ {\alpha}).\tag{5}
\]

为使 \(\pi\) 是 \(\mu\) -可测的，必须且只须 \(\pi\) 是 \(\lambda_{\alpha}\) -可测的，对一切 \(\alpha \in A\) （§1, No.~4, 命题 11 的推论 2）。设这一条件满足；说 \(\pi\) 是 \(\mu\) -逆紧的这时等价于说，对每个函数 \(f \in \mathcal{K}_{+}(X)\) ，

\[
\mu^ {\bullet} (f \circ \pi) <   + \infty .
\]

现在，

\[
\int^ {\bullet} (f \circ \pi) d \mu = \sup _ {\alpha} \int^ {\bullet} (f \circ \pi) d \lambda_ {\alpha} = \sup _ {\alpha} \int^ {\bullet} f d (\pi (\lambda_ {\alpha}))
\]

（§1, No.~4, 命题 11）；因此第一个成员对每个 \(f \in \mathcal{K}_{+}(X)\) 有限，当且仅当族 \((\pi(\lambda_{\alpha}))\) 有一个上确界 \(\theta\) （在 \(\mathcal{M}(X)\) 中），此时 \(\int(f \circ \pi) d\mu = \int f d\theta\) ，这是一个等价于 (5) 的关系式。

推论 1. --- 设 \((\mu_{\alpha})_{\alpha \in A}\) 是 T 上的一个可和的正测度族，使得 \(\mu = \sum_{\alpha \in A} \mu_{\alpha}\) ；为使一个由 T 到局部紧空间 X 中的映射 \(\pi\) 是 \(\mu\) -逆紧的，必须且只须它是 \(\mu_{\alpha}\) -逆紧的，对每个 \(\alpha \in A\) ，且族 \((\pi(\mu_{\alpha}))_{\alpha \in A}\) 是可和的。此时，

\[
\pi (\mu) = \sum_ {\alpha \in A} \pi (\mu_ {\alpha}).\tag{6}
\]

推论 2. --- 设 T 和 X 是两个局部紧空间，\((\lambda_{i})_{1\leqslant i\leqslant n}\) 是 T 上的一个有限正测度序列，并令 \(\mu=\sum_{i=1}^{n}\lambda_{i}\) 。为使一个由 T 到 X 中的映射 \(\pi\) 是 \(\mu\) -逆紧的，必须且只须它对每个指标 i 是 \(\lambda_{i}\) -逆紧的，此时

\[
\sum_ {i = 1} ^ {n} \pi (\lambda_ {i}) = \pi \left(\sum_ {i = 1} ^ {n} \lambda_ {i}\right).
\]

